% !TEX root = chapter-standalone.tex

\section{Traced symmetric monoidal categories}


\todotext{make a remark that without symmetry we would have two distinct concepts of left-trace and right-trace}

\begin{figure}[h!]
    \centering
    \includegraphics[width=10cm]{trace-naturality-x}
    \caption{Naturality in $\Obja$}
    \label{fig:trace-naturality-x}
\end{figure}
\begin{figure}[h!]
    \centering
    \includegraphics[width=10cm]{trace-naturality-y}
    \caption{Naturality in $\Objb$}
    \label{fig:trace-naturality-y}
\end{figure}
\begin{figure}[h!]
    \centering
    \includegraphics[width=10cm]{trace-dinaturality-z}
    \caption{Dinaturality in $\Objc$}
    \label{fig:trace-dinaturality-z}
\end{figure}
\begin{figure}[h!]
    \centering
    \includegraphics[width=6cm]{trace-vanishing-i}
    \caption{Vanishing I}
    \label{fig:trace-vanishing-i}
\end{figure}
\begin{figure}[h!]
    \centering
    \includegraphics[width=10cm]{trace-vanishing-ii}
    \caption{Vanishing II}
    \label{fig:trace-vanishing-ii}
\end{figure}
\begin{figure}[h!]
    \centering
    \includegraphics[width=8cm]{trace-yanking}
    \caption{Yanking}
    \label{fig:trace-yanking}
\end{figure}

\begin{ctdefinition}[Traced monoidal category]
    \label{def:traced-monoidal-cat}
    \label{def:traced-monoidal-category}
    We say that a \SY{symmetric monoidal category}~$\tup{\CatC, \mtimescat, \idmoncat, \associator, \leftunitor, \rightunitor, \braiding}$ is \emph{traced} if it is equipped with a \emph{trace operator}: a family of functions
    \begin{equation}
        \Tr_{\Obja,\Objb}^\Objc\colon \HomSet{\CatC}{\Obja \mtimescatob \Objc}{\Objb\mtimescatob \Objc}\to \HomSet{\CatC}{\Obja}{\Objb},
    \end{equation}
    satisfying the following axioms:
    \begin{enumerate}
        \item \emph{Naturality in $\Obja$:} For any morphisms $\mora\colon \Obja\mtimescatob \Objc \mto \Objb\mtimescatob \Objc$ and $\morb \colon \Obja' \mto \Obja$,
              \begin{equation}\label{eq:trace-naturality-1}
                  \Tr_{\Obja', \Objb}^\Objc ( (\morb \mtimescatmor \catidat\Objc) \mthen \mora) = \morb \mthen \Tr_{\Obja, \Objb}^\Objc (\mora)
              \end{equation}
        \item \emph{Naturality in $\Objb$:} For any morphisms $\mora\colon \Obja\mtimescatob \Objc \mto \Objb\mtimescatob \Objc$ and $\morb \colon \Objb \mto \Objb'$,
              \begin{equation}\label{eq:trace-naturality-2}
                  \Tr_{\Obja, \Objb'}^\Objc ( \mora \mthen (\morb \mtimescatmor \catidat\Objc) ) = \Tr_{\Obja, \Objb}^\Objc (\mora) \mthen \morb
              \end{equation}
        \item \emph{Dinaturality in $\Objc$:}
              For any morphisms $\mora \colon \Obja  \mtimescatob \Objc \mto \Objb \mtimescatob \Objc'$ and $\morb \colon \Objc' \mto \Objc$,
              \begin{equation}\label{eq:trace-dinaturality}
                  \Tr_{\Obja, \Objb}^{\Objc} (\mora \mthen (\catidat\Objb \mtimescatmor \morb)) = \Tr_{\Obja, \Objb}^{\Objc'}((\catidat\Obja \mtimescatmor \morb) \mthen \mora).
              \end{equation}

        \item \emph{Vanishing I:}
              For any morphisms~$\mora \colon \Obja \mtimescatob \idmoncat \mto \Objb \mtimescatob \idmoncat$ in \CatC,
              \begin{equation}
                  \label{eq:vanishing_1}
                  \Tr_{\Obja,\Objb}^\idmoncat (\mora)= \rightunitor^{-1}_{\Obja} \mthen \mora \mthen \rightunitor_\Objb.
              \end{equation}
        \item \emph{Vanishing II:}
              For any morphism~$\mora \colon (\Obja\mtimescatob \Objc) \mtimescatob \Objd \mto (\Objb\mtimescatob \Objc) \mtimescatob \Objd$ in \CatC,
              \begin{equation}
                  \label{eq:vanishing_2}
                  \Tr_{\Obja,\Objb}^\Objc\pars{
                      \Tr_{\Obja \mtimescatob \Objc, \Objb \mtimescatob \Objc}^\Objd(\mora)} = \Tr_{\Obja,\Objb}^{\Objc\mtimescatob \Objd}(\associator_{\Obja, \Objc, \Objd} \mthen \mora \mthen \associator^{-1}_{\Objb, \Objc, \Objd}).
              \end{equation}
        \item \emph{Superposing:}
              For any morphism~$\mora\colon \Obja\mtimescatob \Objc \mto \Objb\mtimescatob \Objc$ in \CatC,
              \begin{equation}
                  \label{eq:superposing}
                  \Tr_{\Obje\mtimescatob \Obja,\Obje\mtimescatob \Objb}^{\Objc}(\associator_{\Obje, \Obja, \Objc} \mthen \catidat\Obje\mtimescatmor \mora \mthen \associator^{-1}_{\Obje, \Objb, \Objc} )=\catidat\Obje\mtimescatmor \Tr_{\Obja,\Objb}^\Objc(\mora).
              \end{equation}
        \item \emph{Yanking:}
              \begin{equation}
                  \label{eq:yanking}
                  \Tr_{\Objc,\Objc}^\Objc\pars{\braiding_{\Objc,\Objc}}=\catidat\Objc.
              \end{equation}
    \end{enumerate}
\end{ctdefinition}

\begin{remark}
    Other variants of the definition of a \SY{traced monoidal category} can be found in the literature.
    For instance, some include a more general version of the superposing law, see \cref{lem:general-superposing-law} below.
\end{remark}

\begin{lemma}
    \label{lem:general-superposing-law}
    Let~$\tup{\CatC, \mtimescat,\idmoncat_{\CatC}, \braiding, \Tr}$ be a \SY{traced monoidal category}.
    Then a more general version of the superposing law holds: for any morphisms~$\mora\colon \Obja\mtimescatob \Objc \mto \Objb \mtimescatob \Objc$ and~$\morb \colon \Objd \mto \Obje$,
    \begin{equation}
        \label{eq:superposing_2}
        \Tr_{\Objd\mtimescatob \Obja,\Obje\mtimescatob \Objb}^{\Objc}(\morb \mtimescatmor \mora)=\morb \mtimescatmor \Tr_{\Obja,\Objb}^\Objc(\mora).
    \end{equation}
\end{lemma}

\todotextjira{165}{\alphubel: @JL: Do the proof.
    Also, this could be made into a graded exercise.
}

\subsection{Relations with cartesian product as monoidal product}

\begin{ctdefinition}
    For the symmetric monoidal category $\tup{\Rel, \cartprod, \singleton, \associator, \leftunitor, \rightunitor, \braiding}$, a trace is defined as follows.
    Given a relation
    \begin{equation}
        \relA \colon \setA \cartprod \setC \to \setB \cartprod \setC,
    \end{equation}
    we set
    \begin{equation}
        \Tr_{\setA,\setB}^\setC (\relA) = \{ \tup{\ela, \elb} \in \setA \cartprod \setB \mid \exists \ \elc \in \setC : \tup{\tup{\ela, \elc}, \tup{\elb, \elc}} \in \relA \}
    \end{equation}
\end{ctdefinition}



\begin{gradedexercise}[\exname{TracingRelations}]

    Given the symmetric monoidal category $\tup{\Rel, \cartprod, \singleton, \associator, \leftunitor, \rightunitor, \braiding}$, consider the trace operation
    $$\Tr_{\Obja,\Objb}^\Objc \colon \HomSet{\RelL}{\Obja \cprod \Objc}{\Objb \cprod \Objc}\to \HomSet{\RelL}{\Obja}{\Objb}$$
    which is defined, for a morphism $\relA \setin\HomSet{\CatC}{\Obja \cprod \Objc}{\Objb\cprod \Objc}$, by
    \begin{equation}
        \Tr_{\Obja,\Objb}^\Objc(\relA) = \makeset{\tup{\ela, \elb} \setin\Obja \cartprod \Objb \mid \exists \elc \setin\Objc \colon \tup{\ela \tupconcat \elc, \elb \tupconcat \elc} \setin\relA}.
    \end{equation}

    Your task is to check that this definition satisfies the following two trace axioms:

    \begin{enumerate}

        \item \emph{Vanishing II}:

              For any relation~$\relA \colon \Obja \cprod \Objc \cprod \Objd \mto \Objb\cprod \Objc \cprod \Objd$ in $\RelL$,
              \begin{equation}
                  \Tr_{\Obja,\Objb}^{\Objc\cprod \Objd}(\relA) = \Tr_{\Obja,\Objb}^\Objc\pars{
                      \Tr_{\Obja \cprod \Objc, \Objb \cprod \Objc}^\Objd(\relA)}.
              \end{equation}

        \item \emph{Superposing}:

              For any relations~$\relA \colon \Obja\cprod \Objc \mto \Objb\cprod \Objc$ and $\relB \colon \Obje \mto \Objf$ in $\RelL$,
              \begin{equation}
                  \Tr_{\Obje\cprod \Obja,\Objf\cprod \Objb}^{\Objc}(\relB \mtimescatmor \relA) = \relB \mtimescatmor \Tr_{\Obja,\Objb}^\Objc(\relA).
              \end{equation}
    \end{enumerate}
\end{gradedexercise}

\solutionof{TracingRelations}


\subsection{Relations with sum as monoidal product}

Consider the symmetric monoidal category $\tup{\Rel, \setdisunion, \Emptyset, \associator, \leftunitor, \rightunitor, \braiding}$.
    In $\Rel$, the disjoint union of any two sets always comes with canonical inclusion relations
    \begin{equation}
        \begin{aligned}
            \inj_1 \colon & \setA \mto \setA\setdisunion\setB \\
                             & \makeset{ \tup{\setAel, \disunionA{\setAel}}},
        \end{aligned}
        \qquad \qquad
        \begin{aligned}
            \inj_2 \colon & \setB  \mto \setA\setdisunion\setB \\
                             & \makeset{ \tup{\setBel, \disunionB{\setBel}}},
        \end{aligned}
    \end{equation}
    and canonical projection relations
    \begin{equation}
        \begin{aligned}
            \proj_1 \colon & \setA\setdisunion\setB \mto \setA \\
                              & \makeset{ \tup{\disunionA{\setAel}, \setAel}},
        \end{aligned}
        \qquad \qquad
        \begin{aligned}
            \proj_2 \colon & \setA\setdisunion\setB  \mto \setB \\
                             & \makeset{ \tup{\disunionA{\setAel}, \setBel}}.
        \end{aligned}
    \end{equation}
    Now suppose we are given a relation of type
    \begin{equation}
        \relA \colon \setA \setdisunion \setC \to \setB \setdisunion \setD.
    \end{equation}
    We can pre-compose it with the inclusion $\inj_1 \colon \setA \sto \setA\setdisunion\setC$ and post-compose it with the projection $\proj_1 \colon \setB \setdisunion\setD \sto \setB$ to obtain a relation
    \begin{equation}
        \relA_{1,1} \colon \setA \overset{\inj_1}{\mto} \setA \setdisunion \setC \overset{\relA}{\mto} \setB \setdisunion \setD \overset{\proj_1}{\mto} \setB.
    \end{equation}
    In a similar manner we also obtain relations
    \begin{equation}
        \relA_{2, 1} \colon \setC \mto \setB, \qquad \relA_{1, 2} \colon \setA \mto \setD, \qquad \relA_{2, 2} \colon \setC \mto \setD,
    \end{equation}
    induced from $\relA$. Together, the four relations obtained in this way contain the same information as the original relation $\relA$, and we can arrange them in a block-matrix notation
    \begin{equation}
\begin{bmatrix}
  \relA_{1,1} & \relA_{2,1} \\
  \relA_{1,2} & \relA_{2,2}
\end{bmatrix}
\end{equation}

which highlights that this decomposition is analogous to block-matrices between vector spaces. 

\begin{ctdefinition}
   For the symmetric monoidal category $\tup{\Rel, \setdisunion, \Emptyset, \associator, \leftunitor, \rightunitor, \braiding}$, a trace is defined as follows.

Given a relation  
   \begin{equation}
       \relA \colon \setA \setdisunion \setC \to \setB \setdisunion \setC,
    \end{equation}
    we set 
    \begin{equation}\label{eq:trace-rel-sum}
        \Tr_{\setA,\setB}^\setC (\relA) = \relA_{1,1} \cup \bigcup_{n \geq 0} ( \relA_{1,2} \mthen \relA_{2,2}^n \mthen \relA_{2,1} ).
    \end{equation}
   Here, ``$\relA_{2,2}^n$'' means ``$\relA_{2,2}$ composed with itself $n$-many times''.
\end{ctdefinition}

\subsection{Co-design}
\todotextjira{201}{\alphubel: @Gioele: Change example}

\begin{marginfigure}
    \centering
    \includesag{50_trace}
    \caption{Design problem with a resource and a functionality of the same type.
    }
    \label{fig:extrace_1}
\end{marginfigure}

\begin{marginfigure}
    \centering
    \includesag{50_trace2}
    \caption{Closing the loop in the design problem.}
    \label{fig:extrace_2}
\end{marginfigure}

Suppose that we are given a design problem with a resource and a functionality of the same type~$\posgenC$ (\cref{fig:extrace_1}).
Can we ``close the loop'', as in the diagram reported in~\cref{fig:extrace_2}?

It turns out that we can give a well-defined semantics to this loop-closing operation, which coincides with the notion of a \emph{trace} in category theory.

The following is the formal definition of the trace operation for \SY{design problems}.

\linkvideo{spring2021-functorial-comp-b:solving-queries:solving-loop} % Loop composition
\begin{ctdefinition}[Trace of a design problem]
    \label{def:dp-trace}
    Given a \SY{design problem}~$\adpa\colon \funposA \mtimescatob \funposC \profto \resposB \mtimescatob \resposC$, its \emph{trace}
    \begin{equation}
        \Tr_{\funposA,\resposB}^\posC(\adpa) \colon \funposA \profto \resposB
    \end{equation}
    is defined as follows:
    %
    \begin{equation}
        \label{eq:tracedef}
        \begin{aligned}
            \Tr_{\funposA,\resposB}^\posC(\adpa) \colon \funposA\posop \Ptimes \resposB & \toinPos \Bool, \\
            \tup{\funposAopel, \RposgenBel}                                             & \mapsto \bigvee_{\posgenCel \setin \posgenC}
            \adpa(\tup{\FposgenAel, \F{\posgenCel}}\Fop,
            \tup{\RposgenBel, \R{\posgenCel}}).
        \end{aligned}
    \end{equation}
\end{ctdefinition}

Think of drawing a loop as a way of writing down the following requirement:
Something that produces~\posC should not use up more of~\posC than it produces.

\begin{lemma}\label{lem:dp-trace-is-trace}
    Trace as in~\cref{def:dp-trace} satisfies the trace axioms.
    In other words,~$\tupp{\DP, {{\otimes}}, \singletonpos, \sigma}$ is a \SY{traced monoidal category}, with trace as in~\cref{eq:tracedef}.
\end{lemma}
\begin{proof}
    We have already shown that~$\tup{\DP,{{\otimes}},\singletonpos,\sigma}$ is a \SY{symmetric monoidal category} (\cref{lem:symmetricmonoidaldp}).
    We prove the trace axioms one by one, starting from vanishing (\cref{eq:vanishing_1}, \cref{eq:vanishing_2}).
    Given any~$\posgenA,\posgenB\setin \Obof\DP$ and~$\adpa\colon \F{\posgenA}\Ptimes \F{\singleton}\profto \R{\posgenB}\Ptimes \R{\singleton}$ in~\DP, we have
    \begin{equation}
        \begin{aligned}
             & = \Tr_{\F{\posgenA},\R{\posgenB}}^{\singleton}(\adpa)(\FposgenAelop,\RposgenBel) \\
             & =\bigvee_{\posCel\setin \singleton}\adpa(\tup{\FposgenAel,\F{\posgenCel}}\Fop, \tup{\RposgenBel,\R{\posgenCel}}) \\
             & =\adpa(\tup{\FposgenAel,\F{\singletonel}}\Fop, \tup{\RposgenBel,\R{\singletonel}}) \\
             & =\adpa(\FposgenAelop,\RposgenBel).
        \end{aligned}
    \end{equation}
    Furthermore, for any morphism~$\adpa\colon \F{\posgenA}\Ptimes \F{\posgenX}\Ptimes \F{\posgenY} \profto \R{\posgenB}\Ptimes \R{\posgenX} \Ptimes \R{\posgenY}$ in \DP, we have
    \begin{equation}
        \begin{aligned}
            ~ & \Tr_{\F{\posgenA},\R{\posgenB}}^{\posX\Ptimes\posY}(\adpa)(\FposgenAelop,\RposgenBel) \\
              & = \bigvee_{\tup{\posXel,\posYel} \setin \posX\Ptimes \posY} \adpa(\tup{\FposgenAel,\F{\posgenXel},\F{\posgenYel}}\Fop, \tup{\RposgenBel,\R{\posgenXel},\R{\posgenYel}}) \\
              & =\bigvee_{\posXel \setin \posX}\pars{\bigvee_{\posYel \setin \posY} \adpa(\tup{\FposgenAel,\F{\posgenXel},\F{\posgenYel}}\Fop, \tup{\RposgenBel,\R{\posgenXel},\R{\posgenYel}})} \\
              & =\Tr_{\F{\posgenA},\R{\posgenB}}^\posX\pars{
                \Tr_{\F{\posgenA}\Ptimes \F{\posgenX},\R{\posgenB}\Ptimes \R{\posgenX}}^\posY(\adpa)(\tup{\FposgenAel,\F{\posgenXel}}\Fop, \tup{\RposgenBel,\R{\posgenXel}})}.
        \end{aligned}
    \end{equation}
    For the superposing axiom (\cref{eq:superposing}), consider~$\adpa\colon \F{\posgenA}\Ptimes \F{\posgenX}\profto \R{\posgenB}\Ptimes \R{\posgenX}$ in \DP.
    We have
    \begin{equation}
        \begin{aligned}
            {}
             & \Tr_{\R{\posgenC}\Ptimes \F{\posgenA},\R{\posgenC}\Ptimes \R{\posgenB}}^{\posX}(\dpidat\posC\mtimescat \adpa)(\tup{\F{\posgenCel_1},\FposgenAel}\Fop, \tup{\R{\posgenCel_2},\RposgenBel}) \\
             & = \bigvee_{\posXel \setin \posX} \dpidat\posC(\F{\posgenCel_1^*},\R{\posgenCel_2})\booland \adpa(\tup{\FposgenAel,\F{\posgenXel}}\Fop, \tup{\RposgenBel,\R{\posgenXel}}) \\
             & =\dpidat\posC(\F{\posgenCel_1^*},\R{\posgenCel_2}) \booland \bigvee_{\posXel\setin \posX} \adpa(\tup{\FposgenAel,\F{\posgenXel}}\Fop, \tup{\RposgenBel,\R{\posgenXel}}) \\
             & =(\dpidat\posC \mtimescat \Tr_{\F{\posgenA},\R{\posgenB}}^\posX(\adpa))(\tup{\F{\posgenCel_1},\FposgenAel}\Fop, \tup{\R{\posgenCel_2},\RposgenBel}).
        \end{aligned}
    \end{equation}
    Finally, for yanking~\cref{eq:yanking} consider~$\sigma_{\posX,\posX}$.
    We have
    \begin{equation}
        \begin{aligned}
            ~ & \Tr_{\posA,\posA}^{\posA}(\sigma_{\posA,\posA})(\F{\posgenAel_1^*},\R{\posgenAel_2}) \\
              & = \bigvee_{\posAel\setin \posA} \sigma_{\posA,\posA}(\tup{\F{\posgenAel_1},\FposgenAel}\Fop, \tup{\R{\posgenAel},\R{\posgenAel_2}}) \\
              & =\bigvee_{\posAel\setin \posA} \F{\posgenAel_1} \posleq \R{\posgenAel_2} \booland \FposgenAel\posleq \R{\posgenAel} \\
              & =\bigvee_{\posAel\setin \posA} \F{\posgenAel_1} \posleq \R{\posgenAel_2} \\
              & =\dpidat\posA(\F{\posgenAel_1^*},\R{\posgenAel_2}).
        \end{aligned}
    \end{equation}
\end{proof}

\begin{gradedexercise}[\exname{DPSnakeTracePart2}]
    \label{ex:DPSnakeTracePart2}

    In this exercise we work again with the category $\DP$ of posets and design problems, equipped with the symmetric monoidal structure where the monoidal product is the cartesian product of posets.
    In the following we make the identification
    \begin{equation}
        (\posgenA \Ptimes \posgenB)\op = \posgenA\op \Ptimes \posgenB\op
    \end{equation}
    for any posets $\posgenA$, $\posgenB$.
    Also, recall that $(\posgenA\op)\op = \posgenA$.

    In components, the associator for $\DP$ is
    \begin{equation}
        \associator_{\posgenA, \posgenB, \posgenC} \colon ((\posgenA \Ptimes \posgenB) \Ptimes \posgenC)\op \Ptimes (\posgenA \Ptimes (\posgenB \Ptimes \posgenC)) \mto \Bool
    \end{equation}
    with
    \begin{equation}
        \associator_{\posgenA, \posgenB, \posgenC}(\tup{\tup{\F{\posgenAel_1\Fop},\F{\posgenBel_1\Fop}}, \F{\posgenCel_1\Fop}},\tup{\R{\posgenAel_2}, \tup{\R{\posgenBel_2}, \R{\posgenCel_2}}}) = \F{\posgenAel_1} \posleq \R{\posgenAel_2} \booland \F{\posgenBel_1} \posleq \R{\posgenBel_2} \booland \F{\posgenCel_1} \posleq \R{\posgenCel_2},
    \end{equation}
    and the left unitor is
    \begin{equation}
        \leftunitor_{\posgenA} \colon (\stylepos{\singleton} \Ptimes \posgenA)\op \Ptimes \posgenA \mto \Bool
    \end{equation}
    with
    \begin{equation}
        \leftunitor_{\posgenA} (\tup{\tup{\F{\bullet\Fop}, \F{\posgenAel_1\Fop}}, \R{\posgenAel_2}}) = \F{\posgenAel_1} \posleq \R{\posgenAel_2}.
    \end{equation}
    The right unitor is analogous.

    The braiding
    $$\braiding_{\posgenA,\posgenB}\colon \F{\posgenA} \Ptimes \F{\posgenB} \profto \R{\posgenB} \Ctimes \R{\posgenA}$$
    is
    \begin{equation}
        \braiding_{\posgenA,\posgenB}(\tup{\F{\posgenAel_1\Fop},\F{\posgenBel_1\Fop}}, \tup{\R{\posgenBel_2},\R{\posgenAel_2}})\definedas \pars{\F{\posgenAel_1}\posleqof\posA \R{\posgenAel_2}}\booland \pars{\F{\posgenBel_1}\posleqof\posB \R{\posgenBel_2}}.
    \end{equation}

    We define the following duality data, with respect to which $\DP$ is compact closed:
    \begin{itemize}
        \item $\posgenA\dual \definedas \posgenA\op$
        \item $\ev_\posgenA \colon \F{(\posgenA\op \Ptimes \posgenA)}\op \Ptimes \R{\makeset{\bullet}} \mto \Bool, \quad \tup{\tup{\ela^*, \elb}^*, \bullet} \mapsto \elb \posleqof\posgenA \ela$
        \item $\coev_\posgenA \colon \F{\makeset{\bullet}}\op \Ptimes \R{(\posgenA \Ptimes \posgenA\op)} \mto \Bool, \quad \tup{\bullet, \tup{\ela, \elb^*}} \mapsto \elb \posleqof\posgenA \ela$
    \end{itemize}

    Your task: given a morphism $\adp \colon \F{\posgenA} \Ptimes \F{\posgenC} \profto \R{\posgenB} \Ctimes \R{\posgenC}$ in $\DP$, show that the design problem $\F{\posgenA} \profto \R{\posgenB}$ given by the following composition
    \begin{equation}
        \rightunitor^{-1}_{\posgenA} \mthen (\catid_\posgenA \mtimescatmor \coev_\posgenC) \mthen \associator^{-1}_{\posgenA, \posgenC, \posgenC\op} \mthen (\adp \mtimescatmor \catid_{\posgenC\op}) \mthen \associator_{\posgenB, \posgenC, \posgenC\op} \mthen (\catid_\posgenB \mtimescatmor \braiding_{\posgenC, \posgenC\op}) \mthen (\catid_\posgenB \mtimescatmor \ev_\posgenC) \mthen \rightunitor_\posgenB
    \end{equation}
    is equal to the design problem $\F{\posgenA} \profto \R{\posgenB}$ given by
    \begin{equation}
        \begin{aligned}
            \funposA\op \Ptimes \resposB    & \toinPos \Bool, \\
            \tup{\funposAopel, \RposgenBel} & \mapsto \bigvee_{\posgenCel \setin \posgenC}
            \adpa(\tup{\FposgenAel, \F{\posgenCel}}\Fop,
            \tup{\RposgenBel, \R{\posgenCel}}).
        \end{aligned}
    \end{equation}
\end{gradedexercise}
\solutionof{DPSnakeTracePart2}

%\begin{widepar}
%    \begin{ctdefinition}[Traced symmetric strict monoidal category]
%        \label{def:traced-sym-strict-mon-cat}
%        We say that a symmetric strict monoidal  \SY{category}~$\tup{\CatC, \mtimescat, \idmoncat, \permmap}$ is \emph{traced} if it is equipped with a family of functions
%        \begin{equation}\label{eq:trace-def-trace-map}
%            \Tr_{\Obja,\Objb}^\Objc\colon \HomSet{\CatC}{\Obja \mtimescatob \Objc}{\Objb\mtimescatob \Objc}\to \HomSet{\CatC}{\Obja}{\Objb},
%        \end{equation}
%        satisfying the following axioms:
%        \begin{enumerate}
%
%            \item \emph{Naturality in $\Obja$:} For any morphisms $\mora\colon \Obja\mtimescatob \Objc \mto \Objb\mtimescatob \Objc$ and $\morb \colon \Obja' \mto \Obja$, if $\Objc$ has an identity $\catidat\Objc$, then
%                  \begin{equation}\label{eq:trace-def-naturality-in-source}
%                      \Tr_{\Obja', \Objb}^\Objc ( (\morb \mtimescatmor \catidat\Objc) \mthen \mora) = \morb \mthen \Tr_{\Obja, \Objb}^\Objc (\mora).
%                  \end{equation}
%
%            \item \emph{Naturality in $\Objb$:}
%                  For any morphisms $\mora\colon \Obja\mtimescatob \Objc \mto \Objb\mtimescatob \Objc$ and $\morb \colon \Objb \mto \Objb'$, if $\Objc$ has an identity $\catidat\Objc$, then
%                  \begin{equation}\label{eq:trace-def-naturality-in-target}
%                      \Tr_{\Obja, \Objb'}^\Objc ( \mora \mthen (\morb \mtimescatmor \catidat\Objc) ) = \Tr_{\Obja, \Objb}^\Objc (\mora) \mthen \morb.
%                  \end{equation}
%
%            \item \emph{Dinaturality in $\Objc$:}
%                  For any morphisms $\mora \colon \Obja \mtimescatob \Objc \mto \Objb \mtimescatob \Objc'$ and $\morb \colon \Objc' \mto \Objc$, if $\catidat\Obja$ and $\catidat\Objb$ exist, then
%                  \begin{equation}\label{eq:trace-def-dinaturality-in-traced-object}
%                      \Tr_{\Obja, \Objb}^{\Objc} (\mora \mthen (\catidat\Objb \mtimescatmor \morb)) = \Tr_{\Obja, \Objb}^{\Objc'}((\catidat\Obja \mtimescatmor \morb) \mthen \mora).
%                  \end{equation}
%
%            \item \emph{Vanishing I:}
%                  For any morphism~$\mora \colon \Obja\mto \Objb$ in \CatC,
%                  \begin{equation}\label{eq:trace-def-vanishing-I}
%                      \Tr_{\Obja,\Objb}^\idmoncat (\mora)=\mora.
%                  \end{equation}
%
%            \item \emph{Vanishing II:}
%                  For any morphism~$\mora \colon \Obja\mtimescatob \Objc \mtimescatob \Objd \mto \Objb\mtimescatob \Objc \mtimescatob \Objd$ in \CatC,
%                  \begin{equation}\label{eq:trace-def-vanishing-II}
%                      \Tr_{\Obja,\Objb}^{\Objc\mtimescatob \Objd}(\mora)=\Tr_{\Obja,\Objb}^\Objc\pars{
%                          \Tr_{\Obja \mtimescatob \Objc, \Objb \mtimescatob \Objc}^\Objd(\mora)}.
%                  \end{equation}
%
%            \item \emph{Superposing:}
%                  For any morphisms~$\mora\colon \Obja\mtimescatob \Objc \mto \Objb\mtimescatob \Objc$ and $\morb \colon \Obje \mto \Objf$ in \CatC,
%                  \begin{equation}\label{eq:trace-def-superposing}
%                      \Tr_{\Obje\mtimescatob \Obja,\Objf\mtimescatob \Objb}^{\Objc}(\morb \mtimescatmor \mora)=\morb \mtimescatmor \Tr_{\Obja,\Objb}^\Objc(\mora).
%                  \end{equation}
%
%            \item \emph{Yanking:}
%                  For any object $\Objc$,
%                  \begin{equation}\label{eq:trace-def-yanking}
%                      \Tr_{\Objc,\Objc}^\Objc(\permmap_{\Objc,\Objc})=\catidat\Objc.
%                  \end{equation}
%        \end{enumerate}
%    \end{ctdefinition}
%\end{widepar}

\todotext{JL: diagram for vanishing II is incorrect -- Z and U need to be swapped}

\subsection{Trace of a linear transformation}
\label{subsec:trace-linear}

Consider the category $\FinVectR$ of finite dimensional real \SY{vector spaces}, which has as objects finite dimensional \SY{vector spaces} and as morphisms linear maps between them.
Using the tensor product $\otimes$ of real \SY{vector spaces} as monoidal product, we can show $\FinVectR$ is a \SY{monoidal category}.
Consider a linear transformation $f\colon B\otimes D\to C\otimes D$, with $B,C,D$ \SY{vector spaces} with bases $\makeset{b_i},\makeset{c_j}$, and $\makeset{d_k}$ respectively.
Here, the trace is a linear function $\mattrace_{B,C}
    ^D(f)\colon B\to C$, given by
\begin{equation}
    \pars{\mattrace_{B,C}
        ^D(f) }_{i,j}=\sum_{k}f_{i\otimes k,j\otimes k}
\end{equation}
\todojira{197}{\alphubel: @JL: Explain notation and show the simple case of sum of diagonals.
    Follow JL's lecture}

\todojira{198}{\alphubel: @Gioele: Add figure using the graphical notations for tensor operations}

\subsection{Trace for symmetric strict monoidal categories}

\begin{ctdefinition}[Trace for $\RelL$ with concatenation as monoidal product]
    \label{def:tuple-REL-trace}
    Consider the symmetric strict monoidal category $\RelL$, where the monoidal product is defined via the concatenation operation $\cprod$ for lists of sets.
    Given a relation $\relA \colon \cObj{\setA, \setC}  \to \cObj{\setB, \setC}$, its \emph{trace} is defined as
    %
    \begin{equation}
        \label{eq:tupel-REL-tracedef}
        \begin{aligned}
            \Tr_{\cObj{\setA}, \cObj{\setB}}^{\cObj{\setC}}(\relA) = \{ \tup{\ela, \elb} \in \cObj{\setA} \times \cObj{\setB} \mid \exists \ \elc \in \cObj{\setC} : \tup{\ela \tupconcat \elc, \elb \tupconcat \elc} \in \relA \}
        \end{aligned}
    \end{equation}
\end{ctdefinition}
